\section{Planes through three points}

Three points can determine a plane, but only if they are not collinear. If the three points lie on one line, then infinitely many planes pass through that line, and the plane is not uniquely determined.

This condition is important. In analytic geometry, we must check whether the data really define the object we want.

\subsection*{From three points to two directions}

Let
\[
A,
\qquad
B,
\qquad
C
\]
be three points in space. If they are not collinear, then the vectors
\[
\overrightarrow{AB}=B-A,
\qquad
\overrightarrow{AC}=C-A
\]
are two non-parallel direction vectors in the plane.

Then the plane can be written parametrically as
\[
P=A+s(B-A)+t(C-A),
\qquad s,t\in\mathbb{R}.
\]

\begin{examplebox}
Find a parametric equation of the plane through
\[
A=(1,0,0),
\qquad
B=(0,1,0),
\qquad
C=(0,0,1).
\]
First compute two direction vectors:
\[
B-A=(-1,1,0),
\]
\[
C-A=(-1,0,1).
\]
Therefore the plane is
\[
P=(1,0,0)+s(-1,1,0)+t(-1,0,1).
\]
In coordinates,
\[
x=1-s-t,
\qquad
y=s,
\qquad
z=t.
\]
\end{examplebox}

This equation says that starting from \(A\), we can move independently toward \(B\) and toward \(C\).

\subsection*{Finding the normal vector}

To obtain the general equation of the plane, we need a normal vector. A normal vector must be perpendicular to both direction vectors in the plane.

If
\[
u=B-A,
\qquad
v=C-A,
\]
then a normal vector is given by the cross product
\[
n=u\times v.
\]
This vector is perpendicular to both \(u\) and \(v\).

\begin{examplebox}
Continue with
\[
A=(1,0,0),
\qquad
B=(0,1,0),
\qquad
C=(0,0,1).
\]
We found
\[
u=(-1,1,0),
\qquad
v=(-1,0,1).
\]
Compute
\[
u\times v=
\begin{vmatrix}
\mathbf{i}&\mathbf{j}&\mathbf{k}\\
-1&1&0\\
-1&0&1
\end{vmatrix}.
\]
Thus
\[
u\times v=(1,1,1).
\]
So a normal vector is
\[
n=(1,1,1).
\]
Using the point \(A=(1,0,0)\), the normal equation is
\[
1(x-1)+1(y-0)+1(z-0)=0.
\]
Therefore
\[
x+y+z-1=0.
\]
Equivalently,
\[
x+y+z=1.
\]
\end{examplebox}

This example is a useful model. Three points give two directions. Two directions give a normal vector. A point and a normal vector give the equation.

\subsection*{Checking non-collinearity}

The points \(A\), \(B\), and \(C\) are collinear exactly when the vectors \(B-A\) and \(C-A\) are parallel. Equivalently, their cross product is the zero vector:
\[
(B-A)\times(C-A)=0.
\]

If the cross product is non-zero, the points determine a unique plane.

\begin{remarkbox}
A zero cross product means that the two direction vectors are not independent. In that case, the three points do not determine a unique plane.
\end{remarkbox}

\begin{summarybox}
To find a plane through three points:
\begin{enumerate}
\item Choose one point as an anchor point, for example \(A\).
\item Compute \(B-A\) and \(C-A\).
\item Check that these vectors are not parallel.
\item Use them as direction vectors for parametric form, or compute their cross product to find a normal vector.
\item Use the normal vector and one point to write the general equation.
\end{enumerate}
\end{summarybox}

The method should be understood as a chain of meanings, not as a memorized procedure.